% RouteSAM 中文双栏理解稿
% 建议使用 XeLaTeX 编译；本稿与英文投稿稿共享公式、图片和参考文献。
\documentclass[UTF8,a4paper,10pt,twocolumn]{ctexart}

\usepackage[left=1.7cm,right=1.7cm,top=1.9cm,bottom=1.9cm]{geometry}
\usepackage{amsmath,amssymb,bm}
\usepackage{graphicx}
\usepackage{float}
\usepackage[authoryear,round]{natbib}
\usepackage[dvipsnames]{xcolor}
\usepackage[colorlinks=true,citecolor=TealBlue,linkcolor=blue,urlcolor=blue]{hyperref}
\usepackage{booktabs}
\usepackage{microtype}

\setlength{\parindent}{2em}
\setlength{\parskip}{0.15em}
\setlength{\columnsep}{0.8cm}
\setlength{\columnseprule}{0pt}
\linespread{1.08}
\raggedbottom

\title{RouteSAM：通过跨图像语义路径实现半监督二维医学图像分割}
\author{作者信息待补充}
\date{}

\begin{document}
\maketitle


\begin{abstract}
医学图像的高昂标注成本使半监督学习成为医学图像分割中的重要研究方向。然而，当有限的标注数据难以充分覆盖多样化的病灶表征时，半监督分割模型往往会出现明显的性能下降。尤其是在未标注图像与现有标注样本之间存在较大语义差异的情况下，有限的标注信息难以稳定迁移至未标注图像，从而影响所生成伪监督的可靠性。
为此，我们提出 RouteSAM，一种基于语义路由的渐进式监督传播方法。为使有限监督覆盖尽可能广泛的训练数据分布，RouteSAM 首先从训练数据中筛选具有较好代表性的样本，作为后续监督传播的起点。进一步地，针对未标注图像，RouteSAM 在其与已标注图像之间构建由多个语义相关图像组成的传播路径，使监督信息沿语义邻近的图像逐步传递，将跨度较大的直接迁移分解为多个更平缓的跨图像传播步骤。
考虑到跨图像传播在不同样本上的可靠性并不一致，我们进一步对传播过程的一致性进行建模，以评估不同伪监督的可靠程度，并据此调节其在半监督训练中的贡献。该策略在保持主路径伪监督不变的同时，降低低可靠传播结果对任务模型学习的干扰，从而更稳定地利用未标注数据扩展有限人工监督。
我们在 ACDC、Kvasir-SEG、BUSI 和 ISIC2018 四个具有不同成像特征与分割目标的医学图像数据集上进行了实验。实验结果表明，RouteSAM 能够在有限标注条件下有效提升未标注数据的利用效率，并在不同医学图像分割任务上取得稳定的性能提升，\textbf{[最终定量结果待补充]}。

\noindent\textbf{关键词：}半监督医学图像分割；Segment Anything Model；跨图像传播；语义路由；伪标签学习
\end{abstract}




\section{引言}

医学图像分割能够对病灶、器官以及其他具有临床意义的解剖结构进行精细的空间勾画，是计算机辅助诊断、疾病定量评估和治疗方案制定等任务的重要基础。近年来，深度学习显著提升了医学图像分割的性能，但高性能分割模型通常需要大量高质量的像素级标注。相比自然图像，医学图像的精确标注往往需要医务人员投入较多时间并具备相应的专业知识，因此大规模医学图像标注数据的获取成本较高且难度较大~\cite{hesamian2019deep,tajbakhsh2020embracing}。为降低对人工标注的依赖，半监督医学图像分割（Semi-Supervised Medical Image Segmentation, SSMIS）通过联合利用少量有标注数据和大量无标注数据，逐渐成为医学图像分割中的重要研究方向~\cite{jiao2024learning}。

现有 SSMIS 方法主要通过伪标签学习与一致性正则化从无标注数据中获取额外监督。一类方法根据模型预测构造伪标签，并通过教师--学生结构或多个网络之间的协同学习不断更新监督信息；另一类方法则对输入、特征表示或模型自身施加扰动，并约束模型在不同扰动条件下保持预测一致。Cross Teaching、MC-Net+ 和 BCP 等代表性方法分别通过异构网络之间的相互监督、不确定性感知的一致性学习，以及有标注与无标注样本之间的双向混合，提高了无标注数据的利用效率~\cite{luo2022crossteaching,wu2022mcnet,bai2023bcp}。然而，当人工标注样本进一步减少时，模型从有标注数据中能够学习到的可靠任务知识也会相应下降，由模型自身产生的伪标签因而更容易受到预测错误和不确定性的影响。

可提示分割基础模型的出现，为获取更强的监督信息带来了新的机会。Segment Anything Model（SAM）及其后续模型通过大规模预训练获得了通用视觉表征能力和提示驱动的分割先验，可以根据点、框或掩码等提示对目标区域进行分割~\cite{kirillov2023segment}。尽管自然图像与医学图像之间存在明显的领域差异，使 SAM 在不同医学影像任务上的直接分割表现并不稳定~\cite{mazurowski2023samexperiment,ma2024medsam}，其通过大规模预训练获得的通用分割知识仍具有较高的利用价值。因此，近年来已有研究开始将 SAM 引入半监督医学图像分割：利用任务特定模型生成提示，并通过 SAM 的预测结果构造或改善无标注数据上的伪监督。例如，SemiSAM 使用任务特定分割模型为 SAM 提供定位提示；SemiSAM+ 进一步建立了通用基础模型与任务特定专用模型之间的协作关系；SAMatch 则借助 SAM 对已有预测结果进行进一步优化~\cite{zhang2023semisam,zhang2025semisamplus,xu2024samatch}。这些研究说明，基础模型中包含的通用分割先验可以成为无标注医学图像学习中的另一种监督来源。然而，现有方法大多仍将重点放在提升每一幅无标注图像自身的预测结果或伪标签质量上，尚未充分利用训练样本之间的跨样本语义关系。

值得注意的是，医学图像之间并非彼此独立，同类解剖结构或病灶在不同样本中往往存在可关联的语义模式。已有研究表明，不同医学图像之间存在可以利用的跨样本语义关系。例如，ACTION 同时建模全局语义关系和局部解剖信息，以提升半监督表示学习效果，这说明训练样本之间并不一定需要被视为彼此完全独立的实例~\cite{you2022action}。进一步地，YoloSeg 使用 SAM2 将单个已标注图像中的分割信息传播到无标注图像，展示了跨图像传播在极低标注环境中的潜力~\cite{zhang2026yoloseg}。然而，这类传播方式仍主要采用从有标注图像直接到目标图像的单跳传播形式。当两幅图像在病灶形态、尺度或外观方面存在较大的语义差异时，直接传播需要在一步之内跨越较大的图像间差异，因而容易产生不稳定的伪监督。这进一步引出了一个问题：是否可以利用无标注数据内部的样本关系，在已有监督与困难目标之间加入若干语义相关的中间图像，从而把一次较困难的直接传播拆解为若干更加局部的传播步骤？

围绕这一问题，我们进一步研究如何利用无标注数据内部的关系来组织有限监督的传播过程。我们将无标注训练集视为由样本间语义关系构成的潜在空间，并提出 \textbf{RouteSeg}。在固定的标注预算下，RouteSeg 首先从训练集中选择能够较广泛覆盖训练数据分布的已标注样本，并将其定义为 \textbf{Anchor}。对于每一个需要生成伪标签的无标注 \textbf{Target}，RouteSeg 从无标注数据中检索若干与之具有语义关联的图像作为 \textbf{Bridge}，再依据图像之间的跨样本关系将其排列成一条有序的\emph{语义路径（semantic route）}：

\begin{equation}
A \rightarrow B_1 \rightarrow B_2 \rightarrow \cdots \rightarrow B_K \rightarrow T.
\end{equation}

随后，RouteSeg 基于 SAM3 沿这条语义路径执行\emph{渐进式跨图像传播（progressive cross-image propagation）}，逐步将 Anchor 中提供的分割信息传播到 Target。相较于直接传播，Bridge 的加入能够将一次跨度较大的图像间变化拆分为多个更加局部的传播过程。这样，无标注样本不再只是等待生成伪标签的 Target，它还可以在其他 Target 的语义路径中充当 Bridge，从而参与有限监督的扩展过程。由于不同的 Anchor 和语义路径可能针对同一个 Target 生成不同的掩码候选，因此直接使用全部传播结果可能会引入不可靠的伪监督。为此，RouteSeg 进一步对这些掩码候选进行质量评估，并选取可靠的\emph{路由伪标签（route-derived pseudo-labels）}用于后续训练。被选中的伪标签与 Anchor 的真实标注共同用于训练任务特定的分割模型，该模型被称为 \textbf{specialist}；SAM3 则作为 \textbf{generalist}。Specialist 从目标医学数据中学习任务特定知识，并进一步对 generalist 的传播结果提供反馈，使基础模型的通用分割先验与医学任务中的领域知识形成互补。语义路由和跨图像传播只在训练阶段使用；在推理阶段，specialist 可以直接对单幅图像完成分割，不需要检索 Anchor、构建语义路径，也不需要额外的人工提示。

我们在 Kvasir-SEG、BUSI 和 ISIC2018 三个具有不同成像特点与病灶类型的二维医学图像分割数据集上评估 RouteSeg，并统一采用约 \(1\%\) 的标注预算。实验将 RouteSeg 与具有代表性的半监督医学图像分割方法、SAM 辅助的低标注分割方法以及直接跨图像传播策略进行比较，从而考察\emph{跨图像监督路由（cross-image supervision routing）}对于扩展有限监督的作用。此外，我们还针对 Anchor 选择、Bridge 的引入和语义路径构建等关键组成部分开展消融与分析，以研究不同路由设计如何影响伪监督质量和最终的分割性能。

本文的主要贡献概括如下：

\begin{itemize}
    \item  我们提出 RouteSeg，一种面向极低标注半监督医学图像分割的跨图像监督路由框架。通过将无标注训练集视为由样本间语义关系构成的潜在空间，RouteSeg 使有限监督能够沿语义路径逐步传播。

    \item 我们提出 Coverage-Aware Anchor Selection 与 Anchor--Bridge--Target 语义路由机制。前者在固定标注预算下选择具有代表性的已标注样本作为 Anchor，后者检索语义相关的无标注 Bridge，并基于 SAM3 执行渐进式跨图像传播，将困难的直接传播拆分为多个局部步骤。

    \item 我们提出可靠的路由伪监督学习与 generalist--specialist 协同策略，从不同语义路径产生的掩码候选中筛选可靠伪标签，并利用 specialist 学习到的任务特定知识进一步改善 generalist 的传播结果。
\end{itemize}


\section{相关工作}

\subsection{半监督医学图像分割}

半监督医学图像分割旨在利用少量像素级标注和大量未标注图像降低医学分割的标注成本。现有方法主要围绕如何从未标注数据中获得可靠监督展开。一类代表性方法通过一致性约束或伪标签学习利用未标注样本。Mean Teacher 使用学生模型参数的指数滑动平均构建教师模型，并约束不同扰动下的预测一致性~\cite{tarvainen2017meanteacher}；Cross Pseudo Supervision 通过两个独立初始化的网络相互生成伪标签~\cite{chen2021cps}；Cross Teaching 进一步利用 CNN 与 Transformer 不同的归纳偏置进行交叉监督~\cite{luo2022crossteaching}；MC-Net+ 则通过多个具有差异性的解码器建模预测不确定性，并利用相互一致性约束未标注数据~\cite{wu2022mcnet}。这些方法逐步从单一教师监督扩展到多模型或异构模型之间的互补监督，其核心仍是从模型自身对未标注图像的预测中构造训练信号。

在此基础上，后续研究进一步关注伪监督的可靠性以及标注数据与未标注数据之间的有效交互。BCP 通过双向 Copy--Paste 在标注样本和未标注样本之间进行区域混合，以减小两类数据之间的经验分布差异~\cite{bai2023bcp}；ABD 根据预测置信度执行自适应的双向区域置换，以控制扰动过程中引入的伪标签噪声~\cite{chi2024abd}；DyCON 则结合动态不确定性感知的一致性学习与对比学习，对不同可靠程度的未标注区域进行差异化优化~\cite{assefa2025dycon}。总体而言，这类方法主要从预测一致性、伪标签质量、数据扰动或局部样本交互的角度提升未标注数据的利用效率。相比之下，本文关注的并不是如何进一步改进单个未标注样本上的预测约束，而是如何利用未标注数据之间的语义关系，显式组织稀疏分割监督在多张图像之间的传播过程。

\subsection{基础模型辅助的少标注医学图像分割}

随着 Segment Anything Model（SAM）等通用分割基础模型的发展，大规模预训练获得的分割先验开始被用于提升少标注医学图像分割中的监督质量。SemiSAM 在传统半监督框架中引入冻结的 SAM 辅助分支，由任务模型根据未标注图像的预测生成提示，并利用 SAM 输出提供额外的一致性监督~\cite{zhang2023semisam}。SemiSAM+ 进一步将这一思路扩展为通用模型与任务模型之间的协同学习，通过置信度感知机制利用一个或多个冻结的基础模型增强未标注数据上的监督信号~\cite{zhang2025semisamplus}。SAMatch 则将弱到强的一致性学习与经过任务适配的 SAM 相结合，由教师模型自动产生提示，再通过 SAM 获得更完整的伪标签用于学生模型训练~\cite{xu2024samatch}。这类方法的共同思路是利用基础模型已有的分割先验，对任务模型产生的定位信息或初始预测进行补充和修正，从而提高单个未标注样本上的伪监督质量。

进一步的研究开始关注基础模型与任务模型之间更充分的双向知识交互。CPC-SAM 通过不同 SAM 解码分支之间的交叉提示与一致性约束增强未标注样本学习~\cite{miao2024cpcsam}；SC-SAM 使用任务专用 U-Net 为 SAM 提供点提示和伪标签以进行参数高效适配，同时利用 SAM 的预测反向约束任务模型~\cite{vu2026scsam}；CPPS-SAM 则通过交叉提示和交叉伪监督，在 CNN 专家模型与 SAM 通用模型之间实现双向知识传递~\cite{zhang2026cppssam}。这些工作表明，通用分割先验与任务特定知识之间的互补能够有效改善少标注条件下的监督学习。然而，现有协同机制主要围绕同一未标注图像上的提示生成、预测修正和模型交互展开，尚未进一步利用不同未标注图像之间的语义联系来组织监督的跨图像传播。RouteSeg 则从这一角度出发，将基础模型的作用由单图伪标签增强进一步扩展到跨图像监督传播过程。

\subsection{跨图像关系建模与监督传播}

除单图像上的一致性学习与伪标签优化外，一些研究开始显式利用不同样本之间的关系来增强半监督表示学习。ACTION 通过全局语义关系与局部解剖对比进行知识蒸馏，以利用数据集内部的跨样本结构~\cite{you2022action}；DACL 构建密度感知的邻域关系，将稀疏区域特征向高密度类别区域聚合~\cite{tang2024dacl}；GraphCL 则通过实例图与图聚类对不同样本间的结构信息进行建模~\cite{wang2024graphcl}。这些工作说明未标注医学图像之间并非相互独立，其潜在的邻域关系和语义结构能够为模型学习提供额外信息。然而，这类方法主要将跨样本关系用于特征约束、对比学习或聚类优化，并未直接将这种关系转化为分割监督在图像之间的传播过程。

另一类与本文更为接近的方法直接利用参考图像或已标注样本向目标图像传递分割信息。PerSAM 根据单个参考图像及其掩码对 SAM 进行个性化，使其能够在目标图像中识别对应对象~\cite{zhang2023persam}；Matcher 通过通用视觉特征匹配实现 one-shot 条件下的目标区域迁移~\cite{liu2023matcher}；OP-SAM 利用支持图像与查询图像之间的特征相关性和迭代提示，将单个标注样本中的息肉先验传递到目标图像~\cite{mao2025opsam}。进一步地，YoloSeg 使用 SAM2 的传播能力将单个标注样本扩展到未标注数据，并通过多视图共识与差异区域学习抑制传播噪声，最终训练任务专用分割模型~\cite{zhang2026yoloseg}。这些方法证明了利用基础模型进行跨图像监督迁移的可行性，但其传播过程主要围绕标注源与目标图像之间的直接对应关系展开。

RouteSeg 进一步关注在标注源与困难目标之间是否存在可用于过渡的未标注样本。与直接从单个支持样本向目标传递监督不同，RouteSeg 将未标注数据集视为一个由样本间语义关系构成的潜在传播空间，并主动选择和排序中间图像，使监督能够沿由 Anchor、Bridge 和 Target 组成的有序路径逐步传递。由此，未标注图像不再仅作为等待生成伪标签的目标样本，而是同时作为连接稀疏标注与困难目标的中间语义节点。





\section{方法}
\label{sec:method}









\subsection{问题定义与整体框架}
\label{sec:framework-overview-zh}

半监督医学图像分割旨在利用少量具有像素级标注的图像与大量未标注图像共同训练分割模型。设训练集为
\begin{equation}
\mathcal{D}_{\mathrm{tr}}
=
\mathcal{D}_{L}\cup\mathcal{D}_{U},
\end{equation}
其中
\begin{equation}
\mathcal{D}_{L}
=
\{(x_i,y_i)\}_{i=1}^{K}
\end{equation}
表示具有人工分割标注的样本集合，
\begin{equation}
\mathcal{D}_{U}
=
\{x_j\}_{j=1}^{N-K}
\end{equation}
表示其余未标注训练图像。当有标注图像仅占整个训练集的极小部分时，少量有标注样本往往只能覆盖有限的病灶形态、尺度和成像外观。当未标注图像与现有有标注样本之间存在较大的视觉或语义差异时，无论是由任务分割网络直接生成伪标签，还是利用可提示分割基础模型进行跨图像迁移，都容易产生不完整、偏移或形态失真的分割结果。更重要的是，一旦这些低质量预测被重新作为训练监督，其误差还可能在后续学习过程中被进一步强化。因此，极低标注条件下的关键问题不仅在于如何为未标注图像产生预测，更在于如何使有限而可靠的监督能够稳定地扩展至更广泛的训练数据。

针对这一问题，我们不再将每一幅未标注图像视为彼此独立的伪标签生成目标，而是利用训练图像之间的语义关系，显式构建监督信息从有标注样本向未标注图像传播的路径。其核心思想是：当某一待处理图像与已有监督样本之间存在较大差异时，不直接在二者之间进行一次跨度较大的监督迁移，而是从训练数据中选择若干语义相关的中间图像，并按照图像之间的相似关系将其组织为从有标注样本逐步到达待处理图像的有序路径。这样，一次困难的直接传播被分解为多个相对平缓的跨图像传播步骤，未标注图像也由单纯的伪标签生成对象进一步转变为能够参与监督传播的中间节点。在此基础上，RouteSAM 按照所构建的语义路径依次向 SAM3 输入图像，并利用其跨帧目标传播能力，使分割信息沿路径逐步传递至待处理图像。对于由此得到的伪监督，我们进一步根据传播过程的一致性估计其可靠程度，并将人工标注与可靠性加权后的伪监督共同用于训练任务特定的分割模型。


整体而言，RouteSAM 主要由四个相互衔接的组成部分构成。首先，为使有限人工监督尽可能覆盖更广泛的训练数据分布，我们从训练集中自动选择一组具有较好分布代表性的样本，作为后续监督传播的起始样本。其次，针对每一幅待处理图像，根据监督起始样本与训练图像之间的语义关系，从训练数据中选择若干语义相关的中间图像，并将其组织为从监督起始样本逐步到达待处理图像的有序语义路径。随后，按照该路径依次向 SAM3 输入图像，利用其跨帧目标传播能力将分割信息逐步传递至待处理图像，从而得到初始伪监督。最后，我们根据传播过程的一致性评估不同伪监督的可靠程度，并据此调节其在半监督训练中的贡献，使人工标注与可靠性加权后的伪监督共同用于训练任务特定的分割模型。上述样本选择、语义路由、跨图像传播以及传播可靠性估计均用于训练阶段的监督扩展；在测试阶段，训练得到的任务特定分割模型直接对单幅医学图像进行预测，无需再次执行监督起始样本选择、语义路由或跨图像传播。








\subsection{覆盖感知的监督起始样本自动选择}
\label{sec:support-selection}

在极低标注条件下，监督起始样本的数量十分有限，其分布直接影响后续监督传播能够覆盖的训练数据范围。当这些样本集中在训练数据分布的局部区域时，部分未标注图像与已有监督之间会存在较大的语义差异，增加跨图像传播过程中的误差累积。针对这一问题，我们在人工标注之前首先从训练集中选择一组具有较好分布代表性的样本，使有限监督尽可能覆盖不同的病灶形态、尺度和成像外观，并将其作为后续语义路由的监督起点。

具体而言，我们使用冻结的 SAM3 视觉编码器提取每幅训练图像的局部特征。SAM3 输出的特征图包含密集的空间 token，相邻位置之间具有较高冗余，因此我们在特征图上均匀采样 64 个位置，并对相应的局部 token 进行 $L_2$ 归一化。这样既保留了图像不同区域的局部信息，也避免直接使用全部密集特征带来的冗余。

随后，将所有训练图像中采样得到的局部特征汇总，并使用 MiniBatch K-Means 构建包含 $M=64$ 个视觉词的共享词典。每个局部特征被分配到距离最近的视觉词，并据此统计每幅图像的视觉词频率。由于部分视觉词可能在大量图像中反复出现，我们采用 TF-IDF 降低这类共有模式的权重。对于图像 $x_i$ 中第 $m$ 个视觉词，其归一化权重定义为
\begin{equation}
p_{im}
=
\frac{
\operatorname{tf}_{im}\operatorname{idf}_{m}
}{
\sum_{l=1}^{M}
\operatorname{tf}_{il}\operatorname{idf}_{l}
},
\qquad
\operatorname{idf}_{m}
=
\log\frac{N+1}{\operatorname{df}_{m}+1}+1,
\end{equation}
其中，$N$ 为训练图像数量，$\operatorname{tf}_{im}$ 表示第 $m$ 个视觉词在图像 $x_i$ 中的出现频率，$\operatorname{df}_{m}$ 表示包含该视觉词的训练图像数量。由此得到图像级视觉词分布
$\mathbf{p}_i=(p_{i1},p_{i2},\ldots,p_{iM})$。

基于视觉词分布，我们使用 Hellinger 映射后的分布相似度比较不同训练图像：
\begin{equation}
k(i,j)
=
\sum_{m=1}^{M}
\sqrt{p_{im}p_{jm}}.
\end{equation}
其中，$k(i,j)$ 越大，表示图像 $x_i$ 与 $x_j$ 的局部视觉模式组成越接近。该度量等价于视觉词分布之间的 Bhattacharyya 系数，使图像间关系由高维局部特征的直接比较转化为局部模式分布的比较。

在此基础上，我们选择一组能够尽可能代表整个训练集的监督起始样本。设所选样本集合为 $\mathcal{A}$，对于任意训练图像 $x_i$，其与集合 $\mathcal{A}$ 中最相近样本的相似度表示该图像被当前集合覆盖的程度。于是定义
\begin{equation}
F(\mathcal{A})
=
\frac{1}{N}
\sum_{i=1}^{N}
\max_{a\in\mathcal{A}}
k(i,a).
\end{equation}
我们通过贪心搜索逐步加入能够最大程度提升 $F(\mathcal{A})$ 的样本，直至达到预设的标注数量。最终选出的样本构成监督起始样本集合，并在后续语义路由和 SAM3 跨图像传播中提供初始监督。











\subsection{面向待处理图像的语义路由构建}
\label{sec:semantic-routing}

尽管上一节选择的监督起始样本能够较好地覆盖训练数据分布，但对于具体的待处理图像而言，其病灶形态、尺度或成像外观仍可能与现有监督起始样本存在较大差异。若直接在两者之间传递分割信息，较大的图像差异容易影响后续传播的稳定性。为此，我们在监督起始样本与待处理图像之间引入若干语义相关的中间图像，将一次跨度较大的跨图像传播分解为多个相对平缓的传播步骤。对于待处理图像 $t$，我们构建有序路径
\begin{equation}
r=(a,x_1,x_2,\ldots,x_6,t),
\end{equation}
其中，$a$ 表示监督起始样本，$x_1,\ldots,x_6$ 表示从训练集中选择的语义中间图像。该路径确定后续监督传播所经过的图像及其顺序。

构建这一路径首先需要判断一幅训练图像是否包含与监督起始样本已标注区域相关的图像内容。由于监督起始样本具有人工分割标注，我们可以利用该标注明确需要向后传播的区域。在 $256\times256$ 分辨率下，我们使用冻结的 SAM3 视觉编码器提取空间局部特征，并将位置 $p$ 处经过 $L_2$ 归一化的特征记为 $\mathbf{z}_{ap}$。随后，将监督起始样本 $a$ 的分割标注映射至对应的特征网格，记标注区域覆盖的位置集合为 $\Omega_a$。对这些位置上的特征进行平均并归一化，得到
\begin{equation}
\mathbf{u}_a
=
\operatorname{norm}
\left(
\frac{1}{|\Omega_a|}
\sum_{p\in\Omega_a}
\mathbf{z}_{ap}
\right).
\end{equation}
$\mathbf{u}_a$ 概括了监督起始样本中人工标注区域的主要局部特征，并用于判断其他训练图像中是否存在与该区域相近的图像内容。

对于任意训练图像 $x_i$，我们同样提取其归一化局部特征 $\{\mathbf{z}_{ip}\}$。为了判断该图像中哪些区域与监督起始样本的已标注区域更为接近，我们首先计算各局部特征与 $\mathbf{u}_a$ 的相似度，并据此进行加权聚合：
\begin{equation}
\alpha_{aip}
=
\frac{
\exp\left(\tau\,\mathbf{u}_a^{\mathrm T}\mathbf{z}_{ip}\right)
}{
\sum_q
\exp\left(\tau\,\mathbf{u}_a^{\mathrm T}\mathbf{z}_{iq}\right)
},
\end{equation}
其中，$\tau$ 为温度系数。由此得到针对监督起始样本 $a$ 的图像表示
\begin{equation}
\mathbf{h}_{a,i}
=
\operatorname{norm}
\left(
\sum_p
\alpha_{aip}\mathbf{z}_{ip}
\right).
\end{equation}
与已标注区域越相近的局部特征在聚合过程中获得越大的权重，因此 $\mathbf{h}_{a,i}$ 更集中地描述了图像 $x_i$ 中与当前监督起始样本相关的部分。进一步定义
\begin{equation}
s(a,i)
=
\mathbf{h}_{a,i}^{\mathrm T}\mathbf{u}_a
\end{equation}
作为图像 $x_i$ 相对于监督起始样本 $a$ 的区域相关性分数。

不同监督起始样本对训练图像产生的相关性分数可能存在不同的整体偏置。例如，某些监督起始样本可能对大量训练图像都产生相对较高的相关性分数，因此仅比较原始 $s(a,i)$ 容易受到这种整体差异的影响。为此，我们以监督起始样本在训练集上的平均相关性为基准，对分数进行中心化：
\begin{equation}
\widetilde{s}(a,i)
=
s(a,i)
-
\frac{1}{|\mathcal{D}_{\mathrm{tr}}|}
\sum_{x_j\in\mathcal{D}_{\mathrm{tr}}}
s(a,j).
\end{equation}
此时，$\widetilde{s}(a,i)$ 表示图像 $x_i$ 相对于训练集平均水平与监督起始样本 $a$ 的已标注区域有多强的相关性。后续路径构建均使用中心化后的 $\widetilde{s}(a,i)$ 作为监督起始样本与路径图像之间的相关性分数。

仅保证路径中的图像与监督起始样本相关仍然不够。两幅图像可能分别与已标注区域具有较高相关性，但二者之间仍存在较大的整体结构或成像外观差异。如果这样的图像被连续排列，路径中仍会出现明显的视觉跳变。因此，我们进一步约束相邻图像之间的整体视觉连续性。对于图像 $x_i$，对其局部特征进行平均并归一化，得到整体描述子
\begin{equation}
\mathbf{g}_i
=
\operatorname{norm}
\left(
\frac{1}{P}
\sum_{p=1}^{P}
\mathbf{z}_{ip}
\right),
\end{equation}
并定义相邻图像 $x_i$ 与 $x_j$ 的视觉相似度为
\begin{equation}
s_{\mathrm{adj}}(i,j)
=
\mathbf{g}_i^{\mathrm T}\mathbf{g}_j.
\end{equation}
因此，一条适合传播的路径需要同时满足两个条件：路径中的图像始终与监督起始样本的已标注区域保持较高相关性，同时相邻图像之间不存在明显的视觉跳变。

基于上述两种关系，我们进一步评价整条候选路径。对于
\begin{equation}
r=(a,x_1,\ldots,x_6,t),
\end{equation}
将路径中各图像相对于监督起始样本的区域相关性以及所有相邻图像之间的视觉相似度汇总为
\begin{equation}
\begin{aligned}
V(r)=&
\{
\widetilde{s}(a,x_1),
\ldots,
\widetilde{s}(a,x_6),
\widetilde{s}(a,t)
\}\\
&\cup
\{
s_{\mathrm{adj}}(a,x_1),
s_{\mathrm{adj}}(x_1,x_2),
\ldots,
s_{\mathrm{adj}}(x_6,t)
\}.
\end{aligned}
\end{equation}
由于一次明显的低相似度跳跃就可能破坏后续传播，我们优先考虑路径中最弱的一环，并在此基础上进一步比较整条路径的平均质量：
\begin{equation}
Q(r)
=
\left(
\min_{v\in V(r)}v,\,
\frac{1}{|V(r)|}
\sum_{v\in V(r)}v
\right).
\end{equation}
不同候选路径按照 $Q(r)$ 进行字典序比较，即优先选择最小值更大的路径；当最小值相同时，再选择平均值更大的路径。这样可以避免路径中某一次明显的跨图像跳变被其他较高的相似度所掩盖。

由于直接枚举所有可能的中间图像组合会带来较大的计算开销，我们采用 beam search 逐步扩展候选路径，并在每一步保留 32 条得分较高的路径。监督起始样本、语义中间图像和待处理图像在同一路径中不重复，所有语义中间图像均来自训练集。

对于待处理图像 $t$，我们进一步在不同路径长度下分别进行搜索。记路径中语义中间图像的数量为 $b$，其中 $b\in\{0,1,\ldots,6\}$，对应的最优路径记为
\[
r_t^{(b)}
=
\left(
a_t^{(b)},
x_1^{(b)},
\ldots,
x_b^{(b)},
t
\right).
\]
当 $b=0$ 时，路径直接由监督起始样本连接至待处理图像；随着 $b$ 增大，路径中逐步引入更多语义中间图像。不同路径长度下的最优路径独立搜索，因此其监督起始样本和中间图像不要求保持一致。

在后续监督传播中，我们以包含六幅语义中间图像的路径 $r_t^{(6)}$ 作为主语义路由，同时保留 $r_t^{(0)},\ldots,r_t^{(5)}$ 作为辅助路径。主路径用于生成实际参与后续半监督训练的伪标签，其余路径则用于评估传播结果对路径变化的稳定性。






\subsection{基于 SAM3 的渐进式跨图像传播}
\label{sec:sam3-propagation}

经过上一节的语义路由构建后，每幅待处理图像 $t$ 均对应 $b=0,\ldots,6$ 七条不同长度的传播路径，其中 $r_t^{(6)}$ 作为主语义路由，其余路径用于后续传播可靠性评估。不同路径长度下的路径独立构建，因此各路径对应的监督起始样本和语义中间图像可以不同。对于每一条路径，我们按照其中确定的图像顺序独立执行一次 SAM3 跨图像传播，将监督起始样本中的分割信息逐步传递至待处理图像。

传播的初始监督来自路径起始处的人工标注。对于路径 $r_t^{(b)}$，监督起始样本 $a_t^{(b)}$ 具有人工分割标注 $Y_{a_t^{(b)}}$。我们根据该标注生成包围标注区域的紧致边界框，并将该边界框作为 SAM3 在路径第一幅图像上的初始提示。除监督起始样本外，路径中的语义中间图像以及待处理图像均不再提供额外的 mask、边界框或文本提示。因此，人工监督仅在路径起点被显式引入，后续图像上的分割结果由 SAM3 沿既定图像顺序逐步传播得到。

具体而言，路径中的所有图像统一调整至 $256\times256$。对于包含 $b$ 幅语义中间图像的路径，其输入顺序为
\begin{equation}
\left[
I_{a_t^{(b)}},
I_{x_1^{(b)}},
\ldots,
I_{x_b^{(b)}},
I_t
\right].
\end{equation}
SAM3 在整个传播过程中保持冻结。模型首先根据监督起始样本上的边界框确定需要传播的分割对象，随后按照语义路由依次处理各中间图像，并利用其跨帧传播能力将分割信息传递至后续图像。与直接从监督起始样本传播至待处理图像相比，引入语义中间图像使监督信息能够按照上一节规划的图像关系逐步到达路径末端。

当传播到待处理图像 $t$ 时，我们将路径 $r_t^{(b)}$ 对应的分割结果记为
\begin{equation}
C_t^{(b)},
\qquad
b\in\{0,1,\ldots,6\}.
\end{equation}
因此，每幅待处理图像最终对应七个传播结果
\begin{equation}
\mathcal{C}_t
=
\left\{
C_t^{(0)},
C_t^{(1)},
\ldots,
C_t^{(6)}
\right\}.
\end{equation}
其中，$C_t^{(0)}$ 来自监督起始样本到待处理图像的直接传播，而 $C_t^{(1)},\ldots,C_t^{(6)}$ 分别对应经过不同数量语义中间图像的传播结果。

在后续半监督学习中，我们固定采用六个语义中间图像路径产生的结果作为待处理图像的初始伪标签，即
\begin{equation}
\widehat{Y}_t
=
C_t^{(6)}.
\end{equation}
其余传播结果 $C_t^{(0)},\ldots,C_t^{(5)}$ 不参与投票、融合或伪标签重构，而是与 $C_t^{(6)}$ 一同保留，用于后续评估伪监督对传播路径变化的稳定性。由此，语义路由不仅为待处理图像提供了实际用于训练的初始伪监督，也为判断该伪监督的传播可靠性提供了多路径预测依据。







\subsection{基于传播可靠性的半监督学习}
\label{sec:reliability-learning}

经过上述跨图像传播，每幅待处理图像都获得了固定的主路径伪标签 $\widehat{Y}_t$，同时保留了其他路径长度对应的传播结果。然而，不同图像上的传播结果具有不同的可靠程度。若在半监督训练中对所有伪标签赋予相同的监督强度，低质量传播结果同样可能被任务模型学习。为此，我们保持主路径产生的伪标签不变，并从传播过程本身估计其可靠程度，据此调节不同伪标签在模型训练中的贡献。

首先，我们利用主语义路由的传播可逆性评估伪标签的闭环稳定性。对于待处理图像 $t$，正向传播沿主路径 $r_t^{(6)}$ 得到伪标签 $\widehat{Y}_t$。随后，在新的 SAM3 传播状态中利用该预测结果重新初始化待处理图像上的分割对象，并沿原路径以相反顺序传播回监督起始样本，得到返回分割结果 $R_t^{(6)}$。由于路径起点具有已知人工标注，因此定义返回一致性为
\begin{equation}
q_{\mathrm{return}}(t)
=
\operatorname{Dice}
\left(
R_t^{(6)},
Y_{a_t^{(6)}}
\right).
\end{equation}

返回一致性衡量的是分割信息经过一次完整的正向和反向传播后，能否重新恢复已知的起始标注。较高的 $q_{\mathrm{return}}(t)$ 表明该传播过程具有较好的闭环稳定性，但并不直接意味着待处理图像上的伪标签一定正确。因此，我们将其作为识别高可靠传播样本的代理信号，而不直接将其作为连续训练权重。具体而言，按照 $q_{\mathrm{return}}(t)$ 从高到低对所有待处理图像进行排序，并将前 $30\%$ 的样本记为高返回一致性集合 $\mathcal{H}$。

对于其余样本，我们进一步利用不同传播路径之间的预测一致性评估伪标签对路径变化的敏感程度。对于同一待处理图像的七个传播结果 $C_t^{(0)},\ldots,C_t^{(6)}$，计算所有两两组合之间的 Dice，并定义多路径一致性为
\begin{equation}
q_{\mathrm{multi}}(t)
=
\frac{1}{21}
\sum_{0\le i<j\le6}
\operatorname{Dice}
\left(
C_t^{(i)},
C_t^{(j)}
\right).
\end{equation}
当不同路径产生的分割结果较为接近时，说明该图像上的预测对路径变化相对稳定；反之，较大的候选差异表明预测结果更依赖具体的路径选择。与返回一致性相同，$q_{\mathrm{multi}}(t)$ 描述的是传播稳定性，而不是伪标签相对于未知真实标注的准确率。

综合上述两种一致性，我们为每幅待处理图像定义图像级伪监督权重
\begin{equation}
w_t=
\begin{cases}
1,
& t\in\mathcal{H},\\[4pt]
0,
& t\notin\mathcal{H}
\ \text{且}\ 
C_t^{(6)}=\varnothing,\\[4pt]
q_{\mathrm{multi}}(t),
& \text{其他情况}.
\end{cases}
\end{equation}
因此，高返回一致性样本采用完整的伪监督权重；对于其余具有非空主路径伪标签的样本，则根据多路径一致性连续调节其监督强度；若主路径未产生有效前景，则该样本不提供伪监督梯度。这里的可靠性权重作用于整幅图像，不进一步引入像素级置信度。

在获得图像级可靠性权重后，我们使用任务特定的 U-Net 联合学习人工标注样本和传播得到的伪标注样本。有标注数据采用较弱的数据增强，伪标注数据采用更强的数据增强。人工标注样本使用交叉熵与 Dice loss 之和作为监督损失 $\mathcal{L}_{\mathrm{sup}}$。对于伪标注样本，同样计算交叉熵与前景 Dice loss，并使用图像级可靠性权重同时调节二者的贡献：
\begin{equation}
\mathcal{L}_{\mathrm{pseudo}}
=
\frac{1}{B_u}
\sum_{t=1}^{B_u}
w_t
\left(
\mathcal{L}_{\mathrm{CE}}^{t}
+
\mathcal{L}_{\mathrm{Dice}}^{t}
\right),
\end{equation}
其中，$B_u$ 表示一个 mini-batch 中的伪标注样本数量。

最终优化目标为
\begin{equation}
\mathcal{L}
=
\mathcal{L}_{\mathrm{sup}}
+
\lambda(k)\mathcal{L}_{\mathrm{pseudo}},
\end{equation}
其中，$\lambda(k)$ 表示随训练进程逐渐增加的伪监督系数，用于在训练初期降低模型对伪标签的依赖，并在模型建立基本分割能力后逐步增加未标注数据的贡献。其具体调度策略与其他优化参数将在实验设置中给出。

整个训练过程中，主语义路由产生的伪标签 $\widehat{Y}_t$ 始终保持不变。返回一致性和多路径一致性仅用于估计其图像级可靠程度，而不参与伪标签投票、融合或重新生成，从而在保留未标注数据监督覆盖的同时降低低可靠传播结果对模型训练的影响。




















\bibliographystyle{plainnat}
\bibliography{../related_work_papers/references}

\end{document}
